\section{Least Action as Survivorship}
\label{sec:least_action}

The Entropic Action defines \textit{what} a system must minimize; the Selection Principle establishes \textit{that} only minimal-action configurations remain observable. We now derive \textit{which} path a
surviving system follows.

Standard quantum mechanics grounds the Principle of Least Action in phase interference. IE suggests a complementary, classical mechanism (thermodynamic elimination) with the substrate itself acting as
a filter: what remains observable is only what has not yet dissipated.

\subsection{The Exponential Filter}
\label{subsec:exponential_filter}

If persistence is maintained by state transitions, and each transition risks scrambling information into the substrate noise floor, then a history of $N$ transitions survives only if every step survives. Independent errors compound multiplicatively, yielding an exponential filter.

Consider a structural history $\Gamma$ composed of $N$ elementary state transitions. Because the substrate possesses a finite noise floor, every transition carries an intrinsic substrate-defined risk of information loss. Let $q$ be the probability that a single transition maintains causal coherence rather than losing its informational trace to the background.

In Quiescent Steady State, successive transitions are statistically independent dissipative events. The probability that the trajectory $\Gamma$ maintains coherence after $N$ transitions is therefore strictly multiplicative:

\begin{equation}
\label{eq:multiplicative_survival}
    P(\text{survival} \mid \Gamma) = q^N
\end{equation}

The Entropic Action $S_\Phi[\Gamma]$ is the dimensionless count of elementary transitions accumulated over the history $\Gamma$. Rewriting $q^N$ in exponential form, $q^N = e^{N\ln q}$, then substituting $N = S_\Phi[\Gamma]$ and $\beta \equiv -\ln q > 0$ (so $\beta$ is positive, since $0<q<1$ makes $\ln q$ negative) gives $N\ln q = -\beta S_\Phi[\Gamma]$, turning \cref{eq:multiplicative_survival} into:

\begin{equation}
\label{eq:survival_prob}
    P(\text{survival} \mid \Gamma) = e^{-\beta S_\Phi[\Gamma]}
\end{equation}

Here $\beta \equiv -\ln q$ is the substrate's inverse entropic tolerance, a geometric invariant of the channel. For discrete substrates, $q$ is determined by the channel geometry; in continuous media, $\beta$ emerges from the substrate's noise spectral density.

The form itself is the informational analogue of the Boltzmann distribution $P \propto e^{-E/k_BT}$, with $S_\Phi$ playing the role of energy and $\beta$ the role of inverse temperature. The Boltzmann form is derived from ensemble statistics at equilibrium and this form is derived from the multiplicative compounding of independent transition risk along a single history, presupposing neither equilibrium, a thermal bath, nor a physical substrate. Both converge on the same exponential law; the informational route recovers it under strictly weaker assumptions.

\subsection{The Severity of the Selection Principle}
\label{subsec:severity}

How harsh is the Selection Principle? Let $\Gamma_{\text{min}}$ be the minimal-action trajectory and $\Gamma_{\text{var}}$ any variant path with a deviation $\delta S_\Phi > 0$. Over any observation window, the relative survival probability is:

\begin{equation}
\label{eq:severity}
    \frac{P(\Gamma_{\text{var}})}{P(\Gamma_{\text{min}})} = e^{-\beta \delta S_\Phi}
\end{equation}

When the fundamental update cycle is much shorter than the observation timescale, even a small per-transition deviation compounds into a vast accumulated $\delta S_\Phi$, and the survival ratio $P(\Gamma_{\text{var}})/P(\Gamma_{\text{min}})$ falls to zero. The gap between the best trajectory and the next-best widens without bound.

\subsection{No Merit Required}
\label{subsec:teleology}
Severity is a statement about the losers, not the winner. It says nothing about whether $\Gamma_{\text{min}}$ is efficient, elegant, or well-suited to anything, only that every alternative to it decayed faster. A mediocre trajectory survives just as certainly as an excellent one, provided everything else was worse. What looks like design, foresight, or purpose is not a property of the winner; it is the absence of its competition.

\subsection{The Principle of Least Action}
\label{subsec:least_action}
The severity of the Selection Principle is formalized as:

\begin{proposition}[Least Action as Survivorship]
In Quiescent Steady State, the observable trajectories of any persistent system are stationary paths of the Entropic Action:
\begin{equation}
    \delta S_\Phi = 0.
\end{equation}
\end{proposition}

\begin{proof}
By the Selection Principle and \cref{eq:severity}, only trajectories minimizing $S_\Phi$ remain observable. The observable ensemble is therefore concentrated on the stationary manifold $\delta S_\Phi = 0$.
\end{proof}

Standard classical mechanics takes the variational principle as foundational. Here, we demonstrate that this principle emerges naturally from a discrete informational substrate as the statistical boundary of survivorship. The record of what is left, not the rule being followed.